\section{Gate contract and verification of the measurements}
\label{app:gates}

Section~\ref{sec:olmo} gives the numerical results for this snapshot. The conditions below specify when a measurement supports the corresponding claim. Each comparison is defined for one checkpoint, one tokenizer and one fixed interface. The exact specification for this edition is given in the \src{observer/gate_contract_ru.md}{contract for the second series}.

\paragraph{Panels and budget.}
Write the eight input bits as $\mathbf{x}=LwR$. Sort the middle strings $w\in\{0,1\}^6$ in ascending order of the SHA-256 digest of the ASCII string \verb!observer-v2|! followed by the six bits, including leading zeros. With ranks starting at one, use 1--8 for development (dev), 9--20 for the joint test of G1 and G2, 21--32 for G3 and 33--44 for G4. Ranks 45--64 form the reserve. For each middle string, present all four endpoint combinations crossed with all four histories, $(L,R,h,e)\in\{0,1\}^4$, in lexicographic order. This gives 16 rows per block, 128 dev rows and 192 rows for each specified 12-block panel. Use dev to select the wording, demonstrations, layer and intervention magnitude. Keep demonstration strings separate from evaluation strings. Freeze the implementations being compared before opening their common held-out panel. A method selected on the basis of results from an opened panel requires a fresh panel.

\paragraph{Obtaining both histories.}
First measure the full-vocabulary probabilities of the single-token continuations \texttt{[space]A} and \texttt{[space]B} at temperature 1. Eligibility requires $P(A),P(B)\ge 0.01$. Use \texttt{random.Random(110)} to draw exactly 1024 samples with replacement from the conditional binary distribution $P(A\mid A\text{ or }B)=P(A)/(P(A)+P(B))$. Record the two original probabilities, the entire sample sequence and the positions of the first A and B. In each history the model processes the first token of the corresponding type that was actually sampled. Take the snapshot after processing this token and before feedback. If either probability falls below the minimum or either letter is absent, assign \texttt{NOT\_EVALUABLE} to this prefix. Forcing a letter or drawing an additional batch of samples would change the experimental condition.

\paragraph{Common criteria.}
For every required numerical row, including neutral, restored and swapped controls, record the target, the full-vocabulary argmax, $\ell_1-\ell_0$ and $q=P(0)+P(1)$. Compute probabilities by normalizing over the full vocabulary in FP32 or higher precision. Each row requires finite values and $q\ge 0.80$. When the margin is strictly positive or negative, the winning bit has probability $>0.40$ and every token outside the pair has probability at most $0.20$. A zero margin counts as a failure and the row remains in the denominator. A complete block requires 16/16 correct outputs. The operating thresholds are 6/8 complete blocks on dev and 9/12 on the held-out panel. Report paired solutions for $L\ne R$ separately for each $e$. The stated 75\% is a fraction of whole blocks. Reusing a prefix does not produce independent samples of model instances.

\paragraph{G0 and G0R.}
G0 checks the weight and tokenizer hashes, the single-token encoding of A/B/X and 0/1 in their actual contexts, exact composition of token IDs, shape, data type and RoPE positions. Diagnostics run on disposable clones. Reading out the past choice stops before the answer is written back. The same primary query must give identical full-vocabulary logits with and without this diagnostic. The original snapshots of both branches retain the bytes of all active K/V entries, the offset and the ordering of visible positions, including ring-buffer indices. The comparison covers the logical data at active positions. Unused allocator capacity outside this range is excluded.

Before selection, G0R reversibly swaps the choice position in all layers, independently of the behavioral score. Given a common preceding history $U$, the snapshots are $C_A=(U,a)$ and $C_B=(U,b)$. For every feedback text and every new string, test both the route $A\to B\to A$ and its mirror starting at B. Apply the reverse swap to the result of the forward swap. Perform readouts on clones of that result. Compare the midpoint with the unmodified donor history and the endpoint with the starting history. Compare metadata, active K/V entries and logits under the same layout. The outputs must be finite, with matching shapes and identical bytes in the original data type. After promoting the logits to FP32, the maximum absolute difference must be zero. Check the reference snapshots before and after. The first Base run used the fixed original \texttt{prefix(8)}. Before swapping, verify the identity of the common prefix and the existence of at least one differing active byte at the choice position. If the branches are already indistinguishable there, G0R returns \texttt{NOT\_EVALUABLE} for that prefix. On dev, two logit comparisons for each of 128 rows give 256 checks.

\paragraph{G1 and G2.}
G1 tests computation from textual feedback on the unmodified histories. In addition to the mandatory numerical probes, use independent copies for D1 (recall $h$ before feedback), D2 (name the corrected rule for the four pairs $h,e$ without the new string) and D3 (apply an explicitly named rule in a fresh episode). Report the letter-valued diagnostics separately from the primary numerical task. Success on the textual control requires the specified format and fraction of complete blocks. For G2, after G0, G0R and G1, repeat the donor and restoration comparisons on panel 9--20. Add a common cache with a neutral token X at the same position and with the same token length. Equality with the unmodified donor follows from $P_B(U,a)=(U,b)$. The positive technical verdict is therefore \texttt{REPLAY\_OK}. G1 already measures behavior on the unmodified histories.

\paragraph{G3.}
After G1 and \texttt{REPLAY\_OK}, replace the textual signal with a single intervention at a neutral token before its K/V entries are written. Specify the token, position, layer, intervention point, vector source and magnitude before testing. Represent a single finite, nonzero $v$ in the operation's data type. Use exact sign reversal to obtain the pulses $v_0=v$, $v_1=-v$. For a given $e$, the two histories receive byte-identical tensors. Check shape, type, elementwise sign reversal and equality of nonzero Euclidean norms computed in FP64 in the same order. For each branch separately, record the norm of the increment after actual rounding, $\operatorname{round}(r+v_e)-r$. Turn off the intervention after this single addition. With zero intervention the two branches of $e$ at fixed $h,\mathbf{x}$ must have identical inputs and logits. Report their accuracy and probability mass. A positive result with the pulse requires the common format and at least 9/12 complete blocks on the G3 panel.

\paragraph{G4.}
Before $\mathbf{x}$ is shown, four donors process the same neutral sequence following $h,e$. Transfer to a fresh template only the K/V entries at prespecified positions and layers of this sequence, denoted by $M_{he}$. All other active positions, tokens, lengths and metadata are identical across the four recipients. The recipients have no direct access to $h,e$ or to any other descendants of the donor history. Define the neutral state $M_*$ once from a common neutral history. The textual channel requires G1 and G2. The pulse channel also requires G3. All of the following criteria must hold on the 12 middle strings:
\begin{enumerate}
\item A recipient given the complete history and the same intervening sequence meets the format criterion and solves at least 9/12 complete blocks. If this control fails, an unsuccessful transfer leaves open whether the donor contained the required information.
\item After inserting $M_*$, all four inputs, active K/V entries and logits at fixed $\mathbf{x}$ are byte-identical. On a single derived snapshot, restoring its own $M_{he}$ after inserting another donor's state restores the original logits.
\item Transferring each donor's own $M_{he}$ meets the format criterion and the threshold of 9/12 complete blocks. Each block includes all four donors and all four endpoint combinations.
\item On the 96 rows with $L\ne R$, let $C_M$ and $C_*$ count correct decisions under transfer and erasure, respectively. Require valid format, $C_*\le48$, 0/12 complete blocks after erasure and $(C_M-C_*)/96\ge 0.25$. Count rows corrected and rows spoiled separately. When one strict binary answer is required, the common input yields exactly 48/96. Nine complete blocks under transfer guarantee at least 72/96 on differing endpoints.
\end{enumerate}

\paragraph{G4-1.}
Before seeing any $\mathbf{x}$, the sender freely selects one token by full-vocabulary greedy decoding. Only the two prespecified single-token completions A/B are admissible. A tie for first place or any other token gives \texttt{FAIL\_FORMAT}, with no retry. Record the four messages in the order of histories $00,01,10,11$. The original Base specification fixes \texttt{prefix(8)}, the literal correct/wrong feedback and a request to record the corrected rule. The measured two-step Instruct condition has its own prompts and an additional threshold $P(A)+P(B)\ge0.80$, fixed in the \src{observer/two-step-interface-04.json}{manifest for variant A}. Treat the two checkpoint conditions separately.

To reproduce the original Base specification, give the sender \verb!\nFeedback: correct! or \verb!\nFeedback: wrong! after its choice, followed by \verb!\nRecord the corrected rule. Rule:!. Here \verb!\n! denotes a single newline. The fresh recipient prefix consists of two lines:
\begin{verbatim}
Rule A selects the first bit. Rule B selects the last bit.
Rule:
\end{verbatim}
Insert the selected ID, which includes a leading space, immediately after the colon. Append \verb!\nBits: !, the eight bits separated by single spaces, and \verb!\nAnswer: ! with its trailing space.

Only the selected ID is delivered to the recipient in the fresh episode. It recomputes the K/V entries for that ID. Check the common-prefix and query-string token IDs against tokenization of the whole text. The external interface excludes the sender's K/V entries and logits, branch index, response time and sampling parameters. Selecting the carrier or branches using $h,e,p$, correct answers or proxies for them is prohibited. For erasure, present the same inputs with a single prespecified constant token A in all worlds. The verdict \texttt{ZERO\_ERROR\_ON\_PANEL} requires four admissible messages, channel-boundary and erasure checks, valid numerical format and 12/12 complete blocks. When testing G4 and G4-1 jointly, freeze both versions before opening panel 33--44 for the first time.

\paragraph{Status precedence and reproducibility.}
A technical error, a nonfinite numerical value or a violation of required byte identity yields \texttt{INVALID}. Next in decreasing order of precedence are \texttt{NOT\_EVALUABLE}, \texttt{FAIL\_FORMAT} and \texttt{FAIL\_DEV}/\texttt{FAIL\_TEST}. A dependent check that was never reached receives \texttt{NOT\_RUN}. A revised interface or a different checkpoint defines a separate condition.

\noindent\begin{minipage}{\linewidth}
Run these two checks from the root of the frozen package.
\begin{verbatim}
python3 -B observer/verify.py
python3 -B observer/verify_part_iii.py
\end{verbatim}
\end{minipage}\par
The first validator recomputes the finite constructions, checks the manifest for the published edition of the historical corpus, verifies the registry addresses and reconstructs the table from the SVG. The second reads the saved dev reports, verifies their condition grids, recorded probabilities and logits, and recomputes accuracy, complete blocks, pairs of opposing histories, diagnostics and the erasure effect. For G0/G0R it rechecks the metadata and the recorded hashes for pairwise comparisons of native logits. Rerunning inference on the GPU and comparing the newly computed tensors byte by byte would constitute a new measurement.
